\documentclass[11pt,a4paper]{article}

\usepackage[margin=2.2cm]{geometry}
\usepackage{amsmath,amssymb,amsthm}
\usepackage{mathtools}
\usepackage{booktabs}
\usepackage{array}
\usepackage{xcolor}
\usepackage{tcolorbox}
\usepackage{enumitem}
\usepackage{fancyhdr}
\usepackage{titlesec}
\usepackage{hyperref}
\usepackage{parskip}
\usepackage{microtype}
\usepackage{multicol}
\usepackage{tabularx}
\usepackage{tikz}

\tcbuselibrary{skins, breakable, theorems}

\definecolor{infoBlue}{RGB}{24,95,165}
\definecolor{infoBg}{RGB}{230,241,251}
\definecolor{successGreen}{RGB}{8,80,65}
\definecolor{successBg}{RGB}{225,245,238}
\definecolor{warnOrange}{RGB}{99,56,6}
\definecolor{warnBg}{RGB}{250,238,218}
\definecolor{dangerRed}{RGB}{121,31,31}
\definecolor{dangerBg}{RGB}{252,235,235}
\definecolor{purpleText}{RGB}{83,74,183}
\definecolor{purpleBg}{RGB}{238,237,254}
\definecolor{tealText}{RGB}{15,110,86}
\definecolor{tealBg}{RGB}{225,245,238}
\definecolor{sectionGray}{RGB}{60,60,60}
\definecolor{coralText}{RGB}{153,60,29}
\definecolor{coralBg}{RGB}{250,236,231}
\definecolor{quizPurple}{RGB}{60,52,137}
\definecolor{quizBg}{RGB}{238,237,254}

\tcbset{
  infobox/.style={colback=infoBg,colframe=infoBlue,left=6pt,right=6pt,top=4pt,bottom=4pt,breakable,boxrule=0.5pt,arc=4pt},
  successbox/.style={colback=successBg,colframe=successGreen,left=6pt,right=6pt,top=4pt,bottom=4pt,breakable,boxrule=0.5pt,arc=4pt},
  warnbox/.style={colback=warnBg,colframe=warnOrange,left=6pt,right=6pt,top=4pt,bottom=4pt,breakable,boxrule=0.5pt,arc=4pt},
  dangerbox/.style={colback=dangerBg,colframe=dangerRed,left=6pt,right=6pt,top=4pt,bottom=4pt,breakable,boxrule=0.5pt,arc=4pt},
  formulabox/.style={colback=gray!8,colframe=gray!35,left=8pt,right=8pt,top=5pt,bottom=5pt,breakable,boxrule=0.5pt,arc=4pt,fontupper=\small},
  questionbox/.style={colback=purpleBg,colframe=purpleText,left=8pt,right=8pt,top=6pt,bottom=6px,boxrule=0.8pt,arc=4pt,fonttitle=\bfseries,coltitle=purpleText,breakable},
  quizbox/.style={colback=quizBg,colframe=quizPurple,left=8pt,right=8pt,top=8pt,bottom=8pt,breakable,boxrule=1pt,arc=6pt,fonttitle=\bfseries\large,coltitle=quizPurple},
  defbox/.style={colback=infoBg,colframe=infoBlue,left=8pt,right=8pt,top=6pt,bottom=6pt,breakable,boxrule=0.5pt,arc=4pt,fonttitle=\bfseries\small,coltitle=infoBlue},
  thmbox/.style={colback=tealBg,colframe=tealText,left=8pt,right=8pt,top=6pt,bottom=6pt,breakable,boxrule=0.5pt,arc=4pt,fonttitle=\bfseries\small,coltitle=tealText},
  exercisebox/.style={colback=gray!5,colframe=gray!40,left=8pt,right=8pt,top=6pt,bottom=6pt,breakable,boxrule=0.6pt,arc=4pt,fonttitle=\bfseries\small,coltitle=sectionGray},
  answerbox/.style={colback=successBg,colframe=successGreen,left=8pt,right=8pt,top=6pt,bottom=6pt,breakable,boxrule=0.5pt,arc=4pt,fonttitle=\bfseries\small,coltitle=successGreen},
  coralbox/.style={colback=coralBg,colframe=coralText,left=6pt,right=6pt,top=4pt,bottom=4pt,breakable,boxrule=0.5pt,arc=4pt}
}

\titleformat{\section}{\large\bfseries\color{sectionGray}}{}{0em}{}[\vspace{-4pt}\rule{\linewidth}{0.4pt}]
\titleformat{\subsection}{\normalsize\bfseries\color{sectionGray}}{\thesubsection.\ }{0em}{}
\titleformat{\subsubsection}{\small\bfseries\itshape\color{sectionGray}}{\thesubsubsection.\ }{0em}{}

\pagestyle{fancy}
\fancyhf{}
\rhead{\small MATH 3032 --- Chapter 4: Hypothesis Testing}
\lhead{\small \textit{Matthew Setiadi \textperiodcentered\ Temple University}}
\rfoot{\small \thepage}
\lfoot{\small \texttt{github.com/mtthwes} \textperiodcentered\ CC BY 4.0}
\renewcommand{\headrulewidth}{0.3pt}

\newcommand{\E}{\mathbb{E}}
\newcommand{\Var}{\operatorname{Var}}
\newcommand{\iid}{\overset{\text{iid}}{\sim}}
\newcommand{\convp}{\xrightarrow{\,p\,}}
\newcommand{\convd}{\xrightarrow{\,d\,}}
\newcommand{\Xbar}{\overline{X}}
\newcommand{\thetahat}{\hat{\theta}_{\text{MLE}}}
\newcommand{\Phifn}{\Phi}

\newcounter{exercise}
\newenvironment{exercise}[1]{%
  \refstepcounter{exercise}%
  \begin{tcolorbox}[exercisebox, title={Exercise \theexercise\ --- #1}]
}{%
  \end{tcolorbox}
}
\newenvironment{solution}{%
  \begin{tcolorbox}[answerbox, title={Solution}]
}{%
  \end{tcolorbox}
}

\newcounter{quizq}
\newenvironment{quizquestion}[1]{%
  \refstepcounter{quizq}%
  \begin{tcolorbox}[quizbox, title={Quiz Question \thequizq\ --- #1}]
}{%
  \end{tcolorbox}
}

\hypersetup{colorlinks=true,linkcolor=infoBlue,urlcolor=infoBlue,
  pdftitle={MATH 3032 Chapter 4 Hypothesis Testing},pdfauthor={Matthew Setiadi}}

% ============================================================
\begin{document}
% ============================================================

\begin{titlepage}
\begin{center}
  \vspace*{1.5cm}
  {\LARGE\bfseries MATH 3032 --- Mathematical Statistics}\\[10pt]
  \rule{\linewidth}{1pt}\\[8pt]
  {\huge\bfseries Chapter 4: Hypothesis Testing}\\[6pt]
  {\Large Deep Understanding Guide + Gaussian Quiz Prep}\\[8pt]
  \rule{\linewidth}{1pt}
  \vspace{1cm}

  \begin{tcolorbox}[formulabox,width=0.85\textwidth,halign=center]
  \centering
  \textbf{Contents of this guide}\\[6pt]
  \begin{tabular}{ll}
    Part I   & The problem setup: why we test hypotheses \\
    Part II  & Error probabilities, significance level, power \\
    Part III & The Neyman-Pearson framework (from lecture slides) \\
    Part IV  & Simple testing and the NP Lemma \\
    Part V   & Unilateral (one-sided) testing \\
    Part VI  & Bilateral (two-sided) testing \\
    Part VII & Unknown variance: the $t$-test \\
    Part VIII & Nuisance parameters \\
    Part IX  & Deep intuition: how to think about all three types \\
    Part X   & Gaussian quiz --- 15 exam-style problems + solutions \\
  \end{tabular}
  \end{tcolorbox}

  \vspace{1cm}
  {\large\bfseries Matthew Setiadi}\\[4pt]
  {\normalsize Accelerated B.S./M.S.\ Computational Data Science}\\[2pt]
  {\normalsize Temple University \textperiodcentered\ \today}

  \vspace{1cm}
  \begin{tcolorbox}[warnbox,width=0.85\textwidth]
  \small\textbf{Quiz focus.} Your quiz is on Gaussian hypothesis testing. Parts VI and VII are the most important. Work every quiz problem in Part X before reading the solution. The quiz will likely require: writing the test function $\delta$, computing the Type I error probability using $\Phi$, computing the Type II error probability, and identifying the rejection region.
  \end{tcolorbox}

  \vfill
  \begin{tcolorbox}[successbox,width=0.85\textwidth]
  \small\centering
  Open-source study notes --- free to use, share, and adapt.\\
  \texttt{github.com/mtthwes} \quad\textperiodcentered\quad Released under CC BY 4.0
  \end{tcolorbox}
\end{center}
\end{titlepage}

\tableofcontents
\newpage

% ============================================================
\section{Part I: Why We Test Hypotheses}
% ============================================================

Chapter 3 answered: \emph{what is the best estimate $\hat\theta$ of $\theta$?} Chapter 4 asks something different: \emph{given data, can we decide whether $\theta$ belongs to one set or another?}

\begin{tcolorbox}[infobox, title={The fundamental question shift}]
\textbf{Estimation (Ch. 3):} Find $\hat\theta$ close to $\theta$.\\
\textbf{Testing (Ch. 4):} Given data, decide: is $\theta \in \Theta_0$ or $\theta \in \Theta_1$?\\[4pt]
These are related but distinct. Testing is about \emph{decision-making under uncertainty}, not just point estimation.
\end{tcolorbox}

\subsection{The Setup (from Lecture Slides)}

\begin{tcolorbox}[defbox, title={Hypothesis test setup}]
Given $X_1,\ldots,X_n \iid F_\theta$, we want to decide between:
\begin{itemize}[leftmargin=1.5em]
  \item $H_0$: the \textbf{null hypothesis}, $\theta \in \Theta_0$ \quad (the ``default'' or ``no effect'' claim)
  \item $H_1$: the \textbf{alternative hypothesis}, $\theta \in \Theta_1$ \quad (the ``signal'' or ``effect'' claim)
\end{itemize}
\textbf{Convention:} Choose $H_0$ as the hypothesis whose false rejection is more harmful.
\end{tcolorbox}

\textbf{Decision tool:} A \textbf{test function} $\delta: \mathcal{X}^n \to \{0,1\}$ where:
\[
  \delta = 1 \;\Leftrightarrow\; \text{reject } H_0 \text{ (decide in favor of } H_1\text{)},\qquad
  \delta = 0 \;\Leftrightarrow\; \text{fail to reject } H_0.
\]

In practice, $\delta$ is built from a \textbf{test statistic} $T(X_1,\ldots,X_n)$ and a \textbf{critical region} $C$:
\[
  \delta(X_1,\ldots,X_n) = \mathbf{1}_{T(X_1,\ldots,X_n)\in C}.
\]

\begin{tcolorbox}[warnbox]
\textbf{$\delta$ is a Bernoulli random variable.} Since it takes values in $\{0,1\}$ and its value depends on the random sample, $\delta \sim \operatorname{Ber}(p)$ where $p = P_\theta(T \in C)$. This $p$ depends on the true $\theta$ --- which is exactly what we're trying to learn about.
\end{tcolorbox}

% ============================================================
\section{Part II: Error Probabilities, Significance, and Power}
% ============================================================

\subsection{The Error Table}

Two mistakes are possible:

\begin{center}
\renewcommand{\arraystretch}{1.5}
\begin{tabular}{@{}lcc@{}}
\toprule
& \textbf{Decide $H_0$ ($\delta=0$)} & \textbf{Decide $H_1$ ($\delta=1$)} \\
\midrule
\textbf{Truth: $H_0$} & No error & \textcolor{dangerRed}{\textbf{Type I error}} \\
\textbf{Truth: $H_1$} & \textcolor{warnOrange}{\textbf{Type II error}} & No error \\
\bottomrule
\end{tabular}
\end{center}

\begin{tcolorbox}[defbox, title={Error probability definitions (Lecture slide 6)}]
\textbf{Type I error probability} (false positive rate):
\[
  h:\Theta_0\to[0,1],\qquad h(\theta) = P_\theta[\delta=1], \quad \theta\in\Theta_0.
\]
\textbf{Type II error probability} (false negative rate):
\[
  g:\Theta_1\to[0,1],\qquad g(\theta) = P_\theta[\delta=0], \quad \theta\in\Theta_1.
\]
\textbf{Power} $= 1-g(\theta) = P_\theta[\delta=1]$ for $\theta\in\Theta_1$. Higher power = better test.
\end{tcolorbox}

\begin{tcolorbox}[dangerbox]
\textbf{Critical warning from lecture:} $h(\theta) \neq 1-g(\theta)$. These functions have \emph{different domains}. $h$ is defined over $\Theta_0$ and $g$ over $\Theta_1$. They are not complements of each other. This is a common source of confusion on exams.
\end{tcolorbox}

\subsection{The Fundamental Trade-off}

From lecture slide 3: if we shrink the critical region $C$ to $C^* \subseteq C$:
\begin{itemize}
  \item Type I error decreases: $P_\theta(T\in C^*) \leq P_\theta(T\in C)$ for $\theta\in\Theta_0$.
  \item Type II error increases: $P_\theta(T\notin C^*) \geq P_\theta(T\notin C)$ for $\theta\in\Theta_1$.
\end{itemize}
\textbf{You cannot minimize both simultaneously.} This is why we need a framework.

% ============================================================
\section{Part III: The Neyman-Pearson Framework}
% ============================================================

\begin{tcolorbox}[thmbox, title={Neyman-Pearson Framework (Lecture slide 4)}]
\textbf{Idea:} Fix the Type I error budget first, then minimize Type II error within that budget.

\begin{enumerate}[leftmargin=2em]
  \item Fix $\alpha \in (0,1)$ --- the \textbf{significance level} (commonly 0.05 or 0.01).
  \item Restrict attention to $\mathcal{D}(\Theta_0,\alpha)$, the class of test functions that \emph{respect the level}:
  \[
    \mathcal{D}(\Theta_0,\alpha) := \left\{\delta : \mathcal{X}^n\to\{0,1\} : \max_{\theta\in\Theta_0} h(\theta) \leq \alpha\right\}.
  \]
  \item Among all $\delta\in\mathcal{D}(\Theta_0,\alpha)$, select the one that \emph{minimizes $g(\theta)$} (equivalently, \emph{maximizes power}).
\end{enumerate}
\end{tcolorbox}

\begin{tcolorbox}[infobox]
\textbf{Intuition:} $\alpha = 0.05$ means ``I am willing to incorrectly reject $H_0$ at most 5\% of the time when $H_0$ is true.'' Given this constraint, find the test that is most likely to correctly reject $H_0$ when $H_1$ is actually true.
\end{tcolorbox}

\subsection{The Three Testing Regimes}

From lecture slide 8, there are three types of hypothesis pairs we study:

\begin{center}
\renewcommand{\arraystretch}{1.6}
\begin{tabular}{@{}llll@{}}
\toprule
\textbf{Type} & \textbf{$H_0$} & \textbf{$H_1$} & \textbf{Key feature}\\
\midrule
Simple & $\theta = \theta_0$ & $\theta = \theta_1$ & Both are single points\\
Unilateral & $\theta \leq \theta_0$ & $\theta > \theta_0$ & One-sided\\
& $\theta \geq \theta_0$ & $\theta < \theta_0$ & One-sided (flipped)\\
Bilateral & $\theta = \theta_0$ & $\theta \neq \theta_0$ & Two-sided\\
\bottomrule
\end{tabular}
\end{center}

\textbf{Deep intuition for which to use:}
\begin{itemize}
  \item \textbf{Simple vs Simple:} You know exactly what both worlds look like. Use the NP Lemma. Optimal test always exists.
  \item \textbf{Unilateral:} You only care about deviations in one direction. E.g., testing if a new drug is \emph{better} (not just different) than a placebo.
  \item \textbf{Bilateral:} You care about deviations in either direction. E.g., testing if a coin is unfair (could be biased toward heads \emph{or} tails).
\end{itemize}

% ============================================================
\section{Part IV: Simple Testing --- The Neyman-Pearson Lemma}
% ============================================================

\subsection{Setup}

$H_0: \theta = \theta_0$ vs $H_1: \theta = \theta_1$ with $\theta_0\neq\theta_1$. Both hypotheses are single points.

\begin{tcolorbox}[thmbox, title={Neyman-Pearson Lemma (Lecture slide 10)}]
Define the \textbf{likelihood ratio}:
\[
  \Lambda(\vec{X}) = \frac{f_{\theta_1}(\vec{X})}{f_{\theta_0}(\vec{X})} = \frac{L(\theta_1)}{L(\theta_0)}.
\]
If there exists $Q>0$ such that $P_{\theta_0}(\Lambda > Q) = \alpha$, then the test
\[
  \delta(\vec{X}) = \mathbf{1}_{\Lambda(\vec{X})>Q}
\]
is \textbf{optimal at level $\alpha$}: it has minimal Type II error probability among all tests at level $\alpha$.
\end{tcolorbox}

\begin{tcolorbox}[infobox]
\textbf{Intuition:} The likelihood ratio $\Lambda$ asks: ``How much more likely is the data under $H_1$ than under $H_0$?'' If $\Lambda$ is large, the data strongly favors $H_1$, so we reject $H_0$. The threshold $Q$ is chosen to keep the Type I error at exactly $\alpha$.
\end{tcolorbox}

\subsection{Worked Example: Exponential Testing (Lecture slides 11--13)}

$X_1,\ldots,X_n \iid \operatorname{Exp}(\lambda)$, test $H_0:\lambda=\lambda_0$ vs $H_1:\lambda=\lambda_1$ with $\lambda_0 < \lambda_1$.

\textbf{Step 1: Likelihood ratio.}
\[
  L(\lambda) = \lambda^n e^{-\lambda\sum X_i}, \qquad
  \Lambda = \frac{L(\lambda_1)}{L(\lambda_0)} = \left(\frac{\lambda_1}{\lambda_0}\right)^n e^{-(\lambda_1-\lambda_0)\sum X_i}.
\]

\textbf{Step 2: Simplify.} $\Lambda$ is a \emph{decreasing} function of $T=\sum X_i$ (since $\lambda_1>\lambda_0$, the exponent is negative times $T$). So:
\[
  \Lambda > Q \iff T \leq q_\alpha \quad \text{for some threshold } q_\alpha.
\]

\textbf{Step 3: Find $q_\alpha$.} Under $H_0$, $T=\sum X_i \sim \operatorname{Gamma}(n,\lambda_0)$. Choose:
\[
  q_\alpha = \alpha\text{-quantile of }\operatorname{Gamma}(n,\lambda_0), \quad\text{i.e.}\quad P_{\lambda_0}(T\leq q_\alpha)=\alpha.
\]

\textbf{Step 4: Test function.}
\[
  \boxed{\delta = \mathbf{1}_{T\leq q_\alpha} = \mathbf{1}_{\sum X_i\leq q_\alpha}}
\]
Reject $H_0$ when the sample sum is unusually \emph{small} (high $\lambda_1$ means shorter waiting times).

\begin{tcolorbox}[successbox]
\textbf{Numerical example from lecture:} $\lambda_0=1$, $\lambda_1=2$, $n=100$, $\alpha=0.05$.\\
$q_\alpha \approx 84.14$ (5th percentile of Gamma$(100,1)$).\\
Reject $H_0$ whenever $\sum X_i < 84.14$, i.e.\ when $\Lambda > 0.000000364$.
\end{tcolorbox}

% ============================================================
\section{Part V: Unilateral (One-Sided) Testing}
% ============================================================

\subsection{The Setup and General Result}

$H_0:\theta\leq\theta_0$ vs $H_1:\theta>\theta_0$ (or flipped). Here $\Theta_0$ is a \emph{set} (a ray), not a single point.

\begin{tcolorbox}[defbox, title={General framework for exponential families (Lecture slide 24)}]
For pdfs in exponential family form $f_\theta(x) = e^{\eta(\theta)K(x)-d(\theta)+S(x)}$, use sufficient statistic $\tau = \sum K(X_i)$.

\begin{center}
\renewcommand{\arraystretch}{1.4}
\begin{tabular}{@{}lll@{}}
\toprule
& $H_0:\theta\leq\theta_0$ vs $H_1:\theta>\theta_0$ & $H_0:\theta\geq\theta_0$ vs $H_1:\theta<\theta_0$\\
\midrule
$\eta(\cdot)$ increasing & $\mathbf{1}\{\tau > q_{1-\alpha}\}$ & $\mathbf{1}\{\tau \leq q_\alpha\}$\\
$\eta(\cdot)$ decreasing & $\mathbf{1}\{\tau \leq q_\alpha\}$ & $\mathbf{1}\{\tau > q_{1-\alpha}\}$\\
\bottomrule
\end{tabular}
\end{center}
\end{tcolorbox}

\subsection{How to Read the Table}

\textbf{Key insight:} Whether the likelihood ratio is increasing or decreasing in the sufficient statistic $\tau$ determines the direction of the rejection region. Here is the logic:

\begin{itemize}
  \item For \textbf{$H_0:\theta\leq\theta_0$ vs $H_1:\theta>\theta_0$}: Evidence against $H_0$ is large $\theta$, which corresponds to large or small $\tau$ depending on whether $\eta(\theta)$ is increasing or decreasing in $\theta$.
  \item \textbf{$\eta$ increasing:} Large $\theta$ $\to$ large $\eta(\theta)$ $\to$ large $\tau$ favors $H_1$ $\to$ reject when $\tau > q_{1-\alpha}$.
  \item \textbf{$\eta$ decreasing:} Large $\theta$ $\to$ small $\eta(\theta)$ $\to$ small $\tau$ favors $H_1$ $\to$ reject when $\tau \leq q_\alpha$.
\end{itemize}

\subsection{Exponential Example (Lecture slides 20--23)}

$X_1,\ldots,X_n \iid \operatorname{Exp}(\lambda)$, test $H_0:\lambda\leq\lambda_0$ vs $H_1:\lambda>\lambda_0$.

The pdf is $f_\lambda(x)=\lambda e^{-\lambda x}$, which in exponential family form is $f_\lambda(x)=e^{\log\lambda+(-\lambda)x}$.

So $\eta(\lambda)=-\lambda$ (decreasing in $\lambda$), $K(x)=x$, sufficient statistic $\tau=\sum X_i$.

Since $\eta$ is \emph{decreasing} and we have $H_0:\theta\leq\theta_0$ vs $H_1:\theta>\theta_0$: reject when $\tau\leq q_\alpha$.

Under $H_0$ evaluated at the boundary $\lambda=\lambda_0$: $T\sim\operatorname{Gamma}(n,\lambda_0)$.
\[
  q_\alpha: \quad P_{\lambda_0}(T\leq q_\alpha) = \alpha \quad\Rightarrow\quad \delta = \mathbf{1}_{T\leq q_\alpha}.
\]

\textbf{Why this is optimal for ALL $\lambda_1>\lambda_0$:} The test $\delta=\mathbf{1}_{T\leq q_\alpha}$ can be rewritten as $\mathbf{1}_{\Lambda>Q(\lambda_1)}$ for any $\lambda_1>\lambda_0$. By the NP Lemma applied pointwise, no other test at level $\alpha$ can beat it for any specific $\lambda_1$. Hence it's \textbf{uniformly most powerful} (UMP).

% ============================================================
\section{Part VI: Bilateral (Two-Sided) Gaussian Testing}
% ============================================================

\subsection{Setup --- Known Variance (Lecture slides 27--30)}

This is the \textbf{most important section for your quiz.}

$X_1,\ldots,X_n \iid \mathcal{N}(\mu,\sigma^2)$ with $\sigma^2$ \textbf{known}. Test $H_0:\mu=\mu_0$ vs $H_1:\mu\neq\mu_0$.

\textbf{Why bilateral?} Unlike unilateral tests, we care about deviations in both directions. If $\mu>\mu_0$, we want to reject; but if $\mu<\mu_0$, we also want to reject.

\subsubsection{Step 1: Form the likelihood ratio}

MLE is $\hat\mu=\overline{X}$. The likelihood ratio is:
\[
  \Lambda = \frac{L(\hat\mu)}{L(\mu_0)} = \exp\!\left(\frac{n}{2\sigma^2}(\overline{X}-\mu_0)^2\right).
\]
Using the algebraic identity (from lecture slide 28):
\[
  \sum(X_i-\mu_0)^2 = \sum(X_i-\overline{X})^2 + n(\overline{X}-\mu_0)^2.
\]

\subsubsection{Step 2: Identify the pivotal quantity}

$\Lambda$ is an increasing function of:
\[
  M := \left(\frac{\overline{X}-\mu_0}{\sigma/\sqrt{n}}\right)^2.
\]
So $\Lambda>Q \iff M>q_{1-\alpha}$ for some threshold $q_{1-\alpha}$.

\subsubsection{Step 3: Identify the distribution under $H_0$}

Under $H_0$, $\mu=\mu_0$, so:
\[
  \overline{X} \sim \mathcal{N}\!\left(\mu_0,\frac{\sigma^2}{n}\right)
  \quad\Rightarrow\quad
  Z := \frac{\overline{X}-\mu_0}{\sigma/\sqrt{n}} \sim \mathcal{N}(0,1)
  \quad\Rightarrow\quad
  M = Z^2 \sim \chi^2(1).
\]

\subsubsection{Step 4: Find the critical value and test function}

\[
  \alpha = P(M>q_{1-\alpha}) = P(\chi^2(1)>q_{1-\alpha}).
\]

Since $\chi^2(1)>c \iff |Z|>\sqrt{c}$ and $P(|Z|>\sqrt{c})=\alpha$:
\[
  2P(Z > z_{1-\alpha/2}) = \alpha \iff P(Z>z_{1-\alpha/2})=\frac{\alpha}{2} \iff z_{1-\alpha/2} = \Phi^{-1}\!\left(1-\frac{\alpha}{2}\right).
\]

\begin{tcolorbox}[thmbox, title={Bilateral Gaussian test (known variance) --- Final form}]
\[
  \boxed{\delta = \mathbf{1}_{\left|\frac{\overline{X}-\mu_0}{\sigma/\sqrt{n}}\right| > z_{1-\alpha/2}}}
\]
where $z_{1-\alpha/2}=\Phi^{-1}(1-\alpha/2)$ is the $(1-\alpha/2)$ quantile of $\mathcal{N}(0,1)$.

\medskip
For $\alpha=0.05$: $z_{0.975}=1.96$. \quad For $\alpha=0.01$: $z_{0.995}=2.576$.
\end{tcolorbox}

\subsection{Type I and Type II Error Probabilities for Gaussian Testing}

\subsubsection{Type I error probability}

The test statistic under $H_0$ is $Z=(\overline{X}-\mu_0)/(\sigma/\sqrt{n})\sim\mathcal{N}(0,1)$.
\[
  h(\mu_0) = P_{\mu_0}(\delta=1) = P(|Z|>z_{1-\alpha/2}) = \alpha. \quad\checkmark
\]

\subsubsection{Type II error probability}

Under $H_1$ with true mean $\mu\neq\mu_0$: $\overline{X}\sim\mathcal{N}(\mu,\sigma^2/n)$, so $Z\sim\mathcal{N}(\mu\sqrt{n}/\sigma - \mu_0\sqrt{n}/\sigma, 1) = \mathcal{N}(\Delta,1)$ where $\Delta = (\mu-\mu_0)\sqrt{n}/\sigma$.

\begin{align*}
  g(\mu) &= P_\mu(\delta=0) = P_\mu(|Z'|<z_{1-\alpha/2}) \quad\text{where } Z'=Z-\Delta\sim\mathcal{N}(0,1)\\
  &= P(-z_{1-\alpha/2}<Z'+\Delta<z_{1-\alpha/2})\\
  &= P(-z_{1-\alpha/2}-\Delta<Z'<z_{1-\alpha/2}-\Delta)\\
  &= \Phi(z_{1-\alpha/2}-\Delta) - \Phi(-z_{1-\alpha/2}-\Delta).
\end{align*}

\begin{tcolorbox}[formulabox]
\textbf{Type II error for bilateral Gaussian test:}
\[
  g(\mu) = \Phi\!\left(z_{1-\alpha/2} - \frac{(\mu-\mu_0)\sqrt{n}}{\sigma}\right) - \Phi\!\left(-z_{1-\alpha/2} - \frac{(\mu-\mu_0)\sqrt{n}}{\sigma}\right)
\]
\textbf{Power} $= 1-g(\mu)$. Power increases as $|\mu-\mu_0|$ increases (easier to detect larger effects) or as $n$ increases (more data).
\end{tcolorbox}

\subsection{The Unilateral Gaussian Tests}

\begin{tcolorbox}[formulabox]
\textbf{Right-sided:} $H_0:\mu\leq\mu_0$ vs $H_1:\mu>\mu_0$
\[
  \delta = \mathbf{1}_{\frac{\overline{X}-\mu_0}{\sigma/\sqrt{n}} > z_{1-\alpha}}, \qquad
  h(\mu_0) = P(Z>z_{1-\alpha}) = \alpha.
\]

\medskip
\textbf{Left-sided:} $H_0:\mu\geq\mu_0$ vs $H_1:\mu<\mu_0$
\[
  \delta = \mathbf{1}_{\frac{\overline{X}-\mu_0}{\sigma/\sqrt{n}} < -z_{1-\alpha}}, \qquad
  h(\mu_0) = P(Z<-z_{1-\alpha}) = \alpha.
\]
\end{tcolorbox}

\textbf{Key differences from bilateral:}
\begin{itemize}
  \item Critical value is $z_{1-\alpha}$ (NOT $z_{1-\alpha/2}$) --- no factor of 2 because rejection is one-directional.
  \item For $\alpha=0.05$: $z_{0.95}=1.645$ (unilateral) vs $z_{0.975}=1.96$ (bilateral).
  \item Unilateral tests are more \emph{powerful} when the true alternative is in the right direction, because the critical value is lower.
\end{itemize}

% ============================================================
\section{Part VII: Unknown Variance --- The Student $t$-Test}
% ============================================================

\subsection{Why We Need a Different Test}

In practice, $\sigma^2$ is rarely known. If we plug in $\hat\sigma^2$ for $\sigma^2$, the statistic $(\overline{X}-\mu_0)/(\hat\sigma/\sqrt{n})$ is no longer $\mathcal{N}(0,1)$ --- it has heavier tails due to estimation uncertainty.

\begin{tcolorbox}[defbox, title={Sample variance and $t$-statistic}]
\textbf{Sample variance} (unbiased):
\[
  S^2 = \frac{1}{n-1}\sum_{i=1}^n(X_i-\overline{X})^2.
\]

\textbf{$t$-statistic} (lecture slide 35):
\[
  T := \frac{\overline{X}-\mu_0}{S/\sqrt{n}}.
\]
\end{tcolorbox}

\begin{tcolorbox}[thmbox, title={Student's $t$-distribution (Lecture slide 36)}]
If $X_1,\ldots,X_n \iid \mathcal{N}(\mu,\sigma^2)$, then under $H_0:\mu=\mu_0$:
\[
  T = \frac{\overline{X}-\mu_0}{S/\sqrt{n}} \sim t_{n-1},
\]
the Student's $t$-distribution with $n-1$ degrees of freedom and pdf:
\[
  f_k(x) = \frac{\Gamma\!\left(\frac{k+1}{2}\right)}{\Gamma\!\left(\frac{k}{2}\right)\sqrt{k\pi}}\left(1+\frac{x^2}{k}\right)^{-\frac{k+1}{2}}, \quad k=n-1.
\]
For $k>2$: $\E[T]=0$ and $\operatorname{Var}(T)=k/(k-2)>1$.
\end{tcolorbox}

\begin{tcolorbox}[thmbox, title={Bilateral $t$-test (Lecture slide 37)}]
$H_0:\mu=\mu_0$ vs $H_1:\mu\neq\mu_0$, unknown $\sigma^2$:
\[
  \boxed{\delta = \mathbf{1}_{\left|T\right|>t_{n-1;\,1-\alpha/2}}}
\]
where $t_{n-1;\,1-\alpha/2}$ is the $(1-\alpha/2)$ quantile of the $t_{n-1}$ distribution.

\medskip
\textbf{Why $|T|>c$?} By symmetry of the $t$-distribution (symmetric around 0, just like $\mathcal{N}(0,1)$), $P(|T|>c)=2P(T>c)=\alpha$ gives $c=t_{n-1;\,1-\alpha/2}$.
\end{tcolorbox}

\subsection{$t$-test vs $Z$-test: When to Use Which}

\begin{center}
\renewcommand{\arraystretch}{1.4}
\begin{tabular}{@{}lll@{}}
\toprule
& \textbf{$Z$-test} & \textbf{$t$-test}\\
\midrule
Variance & Known: $\sigma^2$ given & Unknown: estimate with $S^2$\\
Test stat & $Z=(\overline{X}-\mu_0)/(\sigma/\sqrt{n})$ & $T=(\overline{X}-\mu_0)/(S/\sqrt{n})$\\
Null dist & $\mathcal{N}(0,1)$ & $t_{n-1}$\\
Critical value & $z_{1-\alpha/2}=\Phi^{-1}(1-\alpha/2)$ & $t_{n-1;\,1-\alpha/2}$\\
Large $n$ behavior & Same & $t_{n-1}\to\mathcal{N}(0,1)$, so $t\approx z$\\
\bottomrule
\end{tabular}
\end{center}

\begin{tcolorbox}[infobox]
\textbf{Practical rule:} The $t$-distribution has heavier tails than $\mathcal{N}(0,1)$, so $t_{n-1;\,1-\alpha/2} > z_{1-\alpha/2}$. This means the $t$-test requires more extreme data to reject $H_0$, reflecting the additional uncertainty from estimating $\sigma$. As $n\to\infty$, $t_{n-1;\,1-\alpha/2}\to z_{1-\alpha/2}$, so for large samples it doesn't matter much.
\end{tcolorbox}

% ============================================================
\section{Part VIII: Nuisance Parameters}
% ============================================================

\begin{tcolorbox}[defbox, title={Nuisance parameter setup (Lecture slide 31)}]
When the model has two parameters $(\theta,\xi)$ but we only care about testing $\theta$, the parameter $\xi$ is called a \textbf{nuisance parameter}.

Modified likelihood ratio (generalized LRT):
\[
  \Lambda = \frac{\max_{\theta\neq\theta_0,\,\xi} L(\theta,\xi)}{\max_\xi L(\theta_0,\xi)} = \frac{L(\hat\theta,\hat\xi)_{\text{MLE}}}{\max_\xi L(\theta_0,\xi)}.
\]
\end{tcolorbox}

The bilateral $t$-test is exactly this setup: $\theta=\mu$ is of interest, $\xi=\sigma^2$ is the nuisance parameter. The denominator maximizes over $\sigma^2$ for fixed $\mu_0$; the numerator uses the full MLE.

% ============================================================
\section{Part IX: Deep Intuition --- How to Think About All Three Types}
% ============================================================

\subsection{The Mental Model}

Think of each test as asking a question about where your data falls relative to $\mu_0$:

\begin{tcolorbox}[formulabox]
\textbf{All tests share the same test statistic:}
\[
  Z = \frac{\overline{X}-\mu_0}{\sigma/\sqrt{n}} \quad\text{(or }T=\frac{\overline{X}-\mu_0}{S/\sqrt{n}}\text{ if }\sigma\text{ unknown)}
\]
The \emph{only} difference is what counts as ``extreme enough to reject'':

\begin{center}
\renewcommand{\arraystretch}{1.5}
\begin{tabular}{@{}lll@{}}
\toprule
\textbf{Test type} & \textbf{Reject $H_0$ when} & \textbf{Intuition}\\
\midrule
Bilateral ($\mu\neq\mu_0$) & $|Z|>z_{1-\alpha/2}$ & Too far from $\mu_0$ in either direction\\
Right-sided ($\mu>\mu_0$) & $Z>z_{1-\alpha}$ & Too far above $\mu_0$\\
Left-sided ($\mu<\mu_0$) & $Z<-z_{1-\alpha}$ & Too far below $\mu_0$\\
\bottomrule
\end{tabular}
\end{center}
\end{tcolorbox}

\subsection{Why $z_{1-\alpha}$ for Unilateral but $z_{1-\alpha/2}$ for Bilateral?}

\textbf{The $\alpha$-budget.} For $\alpha=0.05$:
\begin{itemize}
  \item \textbf{Bilateral:} split the budget equally between the two tails: $\alpha/2=0.025$ in each tail. So $z_{0.975}=1.96$.
  \item \textbf{Unilateral:} put all the budget in one tail: $\alpha=0.05$ in one tail. So $z_{0.95}=1.645$.
\end{itemize}
Result: $1.645 < 1.96$, so right-sided tests reject more easily (lower bar) when $\mu>\mu_0$, but cannot detect $\mu<\mu_0$ at all.

\subsection{The $\Phi$ function --- Your Main Calculation Tool}

$\Phi(x) = P(Z\leq x)$ for $Z\sim\mathcal{N}(0,1)$. Key properties:
\begin{itemize}
  \item $\Phi(-x) = 1-\Phi(x)$ (symmetry)
  \item $P(Z>x) = 1-\Phi(x)$
  \item $P(|Z|>x) = 2(1-\Phi(x)) = 2\Phi(-x)$
  \item $P(a<Z<b) = \Phi(b)-\Phi(a)$
\end{itemize}

\begin{tcolorbox}[formulabox]
\textbf{Critical values you must memorize:}
\begin{center}
\begin{tabular}{@{}cccc@{}}
\toprule
$\alpha$ & $z_{1-\alpha}$ (unilateral) & $z_{1-\alpha/2}$ (bilateral) & $t_{n-1;1-\alpha/2}$ (large $n$)\\
\midrule
0.10 & 1.282 & 1.645 & $\approx 1.645$\\
0.05 & 1.645 & 1.960 & $\approx 1.960$\\
0.01 & 2.326 & 2.576 & $\approx 2.576$\\
\bottomrule
\end{tabular}
\end{center}
\end{tcolorbox}

% ============================================================
\section{Part X: Gaussian Hypothesis Testing Quiz --- 15 Problems}
% ============================================================

\begin{tcolorbox}[quizbox, title={Quiz instructions}]
These 15 problems cover everything your quiz will test. Work each problem completely before reading the solution. For Gaussian problems: (1) identify $H_0$, $H_1$, and test type; (2) write $Z$ or $T$; (3) state the rejection region; (4) compute Type I error in terms of $\Phi$; (5) compute Type II error if asked.

\textbf{Notation:} $\Phi$ = standard normal CDF. $z_p = \Phi^{-1}(p)$ is the $p$-quantile.
\end{tcolorbox}

\subsection{Warm-Up: Understanding the Setup}

\begin{quizquestion}{Identifying the test type}
For each pair $(H_0, H_1)$, state whether it is Simple, Right-sided Unilateral, Left-sided Unilateral, or Bilateral:
\begin{enumerate}[label=(\alph*)]
  \item $H_0:\mu=5$ vs $H_1:\mu=7$
  \item $H_0:\mu=0$ vs $H_1:\mu\neq 0$
  \item $H_0:\mu\leq 3$ vs $H_1:\mu>3$
  \item $H_0:\mu\geq 10$ vs $H_1:\mu<10$
  \item $H_0:\sigma^2=4$ vs $H_1:\sigma^2\neq 4$ \quad (variance test, not mean)
\end{enumerate}
\end{quizquestion}

\begin{solution}
(a) Simple (both hypotheses are single points, $\theta_0=5$, $\theta_1=7$).\\
(b) Bilateral ($H_1$ is $\mu\neq 0$, two-sided).\\
(c) Right-sided unilateral ($H_1:\mu>3$).\\
(d) Left-sided unilateral ($H_1:\mu<10$).\\
(e) Bilateral for variance (same structure as bilateral mean test but for $\sigma^2$).
\end{solution}

\begin{quizquestion}{Writing the test function}
$X_1,\ldots,X_{25} \iid \mathcal{N}(\mu,16)$ (so $\sigma^2=16$, $\sigma=4$). Test $H_0:\mu=10$ vs $H_1:\mu\neq 10$ at level $\alpha=0.05$.
\begin{enumerate}[label=(\alph*)]
  \item Write the test statistic $Z$ explicitly.
  \item State the rejection region.
  \item Write the test function $\delta$.
  \item What is the Type I error probability?
\end{enumerate}
\end{quizquestion}

\begin{solution}
\textbf{(a)} $Z = \dfrac{\overline{X}-10}{4/\sqrt{25}} = \dfrac{\overline{X}-10}{4/5} = \dfrac{5(\overline{X}-10)}{4}$.

\textbf{(b)} Bilateral test: reject when $|Z|>z_{0.975}=1.96$.

\textbf{(c)} $\delta = \mathbf{1}_{|Z|>1.96} = \mathbf{1}_{\left|\frac{5(\overline{X}-10)}{4}\right|>1.96}$.

\textbf{(d)} $h(\mu_0) = P_{\mu_0}(|Z|>1.96) = 2(1-\Phi(1.96)) = 2(0.025) = 0.05 = \alpha$. \checkmark
\end{solution}

\begin{quizquestion}{Type I error computation}
$X_1,\ldots,X_n \iid \mathcal{N}(\mu,1)$. Test $H_0:\mu=0$ vs $H_1:\mu\neq 0$ using test $\delta=\mathbf{1}_{|\overline{X}|\geq Q}$.
\begin{enumerate}[label=(\alph*)]
  \item Express $h(0)=P_{\mu=0}(\delta=1)$ in terms of $\Phi$, $Q$, $n$.
  \item For $\alpha=0.05$, find the smallest $Q$ such that $h(0)\leq 0.05$.
\end{enumerate}
This is exactly the computation from the lecture slides (Introduction to Hypothesis Testing, slides 8--9).
\end{quizquestion}

\begin{solution}
\textbf{(a)} Under $H_0$: $\overline{X}\sim\mathcal{N}(0,1/n)$. Standardize: $U=\overline{X}/(1/\sqrt{n})=\sqrt{n}\,\overline{X}\sim\mathcal{N}(0,1)$.
\[
h(0) = P_0(|\overline{X}|\geq Q) = P(|\overline{X}|/(1/\sqrt{n})\geq Q\sqrt{n}) = P(|U|\geq Q\sqrt{n}) = 2\Phi(-Q\sqrt{n}).
\]

\textbf{(b)} Require $2\Phi(-Q\sqrt{n})\leq 0.05$, i.e.\ $\Phi(-Q\sqrt{n})\leq 0.025$, i.e.\ $-Q\sqrt{n}\leq\Phi^{-1}(0.025)=-1.96$.
So $Q\geq 1.96/\sqrt{n}$. The smallest such $Q$ is $Q=1.96/\sqrt{n}$.
\end{solution}

\subsection{Core Calculations}

\begin{quizquestion}{Type II error for bilateral test}
$X_1,\ldots,X_{100}\iid\mathcal{N}(\mu,\sigma^2)$ with $\sigma=2$. Test $H_0:\mu=5$ vs $H_1:\mu\neq 5$ at $\alpha=0.05$.
\begin{enumerate}[label=(\alph*)]
  \item Write the test function $\delta$.
  \item Compute $g(6) = P_{\mu=6}(\delta=0)$ (Type II error when true mean is 6).
  \item Compute the power at $\mu=6$.
\end{enumerate}
\end{quizquestion}

\begin{solution}
\textbf{(a)} $Z=(\overline{X}-5)/(2/\sqrt{100})=(\overline{X}-5)/0.2=5(\overline{X}-5)$.
$\delta = \mathbf{1}_{|Z|>1.96}$.

\textbf{(b)} Under $\mu=6$: $\overline{X}\sim\mathcal{N}(6, 4/100)$, so $\overline{X}\sim\mathcal{N}(6,0.04)$.
Write $Z=5(\overline{X}-5)$; under $\mu=6$: $Z\sim\mathcal{N}(5,1)$ since $5(\mu-5)=5(1)=5=\Delta$.

$g(6) = P_6(|Z|\leq 1.96) = P_6(-1.96\leq Z\leq 1.96)$.
Let $Z'=Z-5\sim\mathcal{N}(0,1)$:
\[
g(6) = P(-1.96-5\leq Z'\leq 1.96-5) = P(-6.96\leq Z'\leq -3.04) = \Phi(-3.04)-\Phi(-6.96).
\]
$\Phi(-3.04)\approx 0.0012$, $\Phi(-6.96)\approx 0$. So $g(6)\approx 0.0012$.

\textbf{(c)} Power $=1-g(6)\approx 0.9988$. Very high --- $\mu=6$ is far from $\mu_0=5$ relative to $\sigma/\sqrt{n}=0.2$.
\end{solution}

\begin{quizquestion}{Right-sided unilateral test}
$X_1,\ldots,X_{36}\iid\mathcal{N}(\mu,9)$. Test $H_0:\mu\leq 0$ vs $H_1:\mu>0$ at $\alpha=0.05$.
\begin{enumerate}[label=(\alph*)]
  \item Write the test function $\delta$.
  \item Verify $\sup_{\mu\leq 0} h(\mu)\leq\alpha$.
  \item Compute the power at $\mu=1$.
\end{enumerate}
\end{quizquestion}

\begin{solution}
\textbf{(a)} $\sigma=3$, $\sigma/\sqrt{n}=3/6=0.5$. $Z=(\overline{X}-0)/0.5=2\overline{X}$.
Right-sided: reject when $Z>z_{0.95}=1.645$.
$\delta = \mathbf{1}_{2\overline{X}>1.645} = \mathbf{1}_{\overline{X}>0.8225}$.

\textbf{(b)} For any $\mu\leq 0$: $Z\sim\mathcal{N}(2\mu,1)$ (mean $\leq 0$).
$h(\mu)=P_\mu(Z>1.645)=P(\mathcal{N}(2\mu,1)>1.645)=1-\Phi(1.645-2\mu)$.
Since $\mu\leq 0$: $1.645-2\mu\geq 1.645$, so $\Phi(1.645-2\mu)\geq\Phi(1.645)=0.95$, giving $h(\mu)\leq 0.05$. $\checkmark$
Supremum achieved at $\mu=0$: $h(0)=1-\Phi(1.645)=0.05=\alpha$.

\textbf{(c)} At $\mu=1$: $Z\sim\mathcal{N}(2,1)$. Power $=P_1(Z>1.645)=P(\mathcal{N}(2,1)>1.645)=1-\Phi(1.645-2)=1-\Phi(-0.355)=\Phi(0.355)\approx 0.639$.
\end{solution}

\begin{quizquestion}{Left-sided test}
$X_1,\ldots,X_{16}\iid\mathcal{N}(\mu,4)$. Test $H_0:\mu\geq 3$ vs $H_1:\mu<3$ at $\alpha=0.01$.
\begin{enumerate}[label=(\alph*)]
  \item Write $\delta$.
  \item What is $h(3)$?
  \item If observed $\overline{x}=2.1$, do you reject $H_0$?
\end{enumerate}
\end{quizquestion}

\begin{solution}
\textbf{(a)} $\sigma/\sqrt{n}=2/4=0.5$. Left-sided: reject when $Z<-z_{0.99}=-2.326$.
$Z=(\overline{X}-3)/0.5=2(\overline{X}-3)$.
$\delta=\mathbf{1}_{2(\overline{X}-3)<-2.326}=\mathbf{1}_{\overline{X}<3-1.163}=\mathbf{1}_{\overline{X}<1.837}$.

\textbf{(b)} $h(3)=P_3(Z<-2.326)=\Phi(-2.326)=0.01=\alpha$. \checkmark

\textbf{(c)} $\overline{x}=2.1$. Is $2.1<1.837$? No. Do not reject $H_0$.
\end{solution}

\begin{quizquestion}{Finding sample size for desired power}
$X_1,\ldots,X_n\iid\mathcal{N}(\mu,1)$. Test $H_0:\mu=0$ vs $H_1:\mu\neq 0$ at $\alpha=0.05$. What is the minimum $n$ such that power at $\mu=0.5$ is at least $0.80$?
\end{quizquestion}

\begin{solution}
Test statistic: $Z=\overline{X}\sqrt{n}\sim\mathcal{N}(0,1)$ under $H_0$; $Z\sim\mathcal{N}(0.5\sqrt{n},1)$ under $\mu=0.5$.

Power at $\mu=0.5$:
\[
1-g(0.5) = P_{\mu=0.5}(|Z|>1.96) \approx P_{\mu=0.5}(Z>1.96) = 1-\Phi(1.96-0.5\sqrt{n}).
\]
(The other tail $P(Z<-1.96)$ under $\mu=0.5$ is negligible.)

Require $1-\Phi(1.96-0.5\sqrt{n})\geq 0.80$, i.e.\ $\Phi(1.96-0.5\sqrt{n})\leq 0.20$.

$\Phi^{-1}(0.20)=-0.842$. So $1.96-0.5\sqrt{n}\leq -0.842$, giving $0.5\sqrt{n}\geq 2.802$, so $\sqrt{n}\geq 5.604$, thus $n\geq 32$ (round up).

\textbf{Answer:} $n=32$.
\end{solution}

\subsection{The $t$-Test}

\begin{quizquestion}{$t$-test setup}
$X_1,\ldots,X_{10}\iid\mathcal{N}(\mu,\sigma^2)$ with $\sigma^2$ \emph{unknown}. Observed $\overline{x}=5.3$, $s^2=4$ (so $s=2$). Test $H_0:\mu=4$ vs $H_1:\mu\neq 4$ at $\alpha=0.05$.
\begin{enumerate}[label=(\alph*)]
  \item Compute the $t$-statistic.
  \item State the rejection region using $t_{9;\,0.975}=2.262$.
  \item Do you reject $H_0$?
\end{enumerate}
\end{quizquestion}

\begin{solution}
\textbf{(a)} $T=(\overline{x}-\mu_0)/(s/\sqrt{n})=(5.3-4)/(2/\sqrt{10})=1.3/(2/3.162)=1.3/0.632\approx 2.057$.

\textbf{(b)} Reject when $|T|>t_{9;\,0.975}=2.262$.

\textbf{(c)} $|T|=2.057<2.262$. Do NOT reject $H_0$. (Just barely within the non-rejection region.)
\end{solution}

\begin{quizquestion}{$t$-test Type I error expression}
$X_1,\ldots,X_n\iid\mathcal{N}(\mu,\sigma^2)$, $\sigma^2$ unknown. Test $H_0:\mu=\mu_0$ vs $H_1:\mu\neq\mu_0$ using the $t$-statistic $T=(\overline{X}-\mu_0)/(S/\sqrt{n})$.
\begin{enumerate}[label=(\alph*)]
  \item What is the distribution of $T$ under $H_0$?
  \item Write the Type I error probability $h(\mu_0)$ in terms of the $t$-distribution CDF.
  \item Why is $t_{n-1;\,1-\alpha/2}>z_{1-\alpha/2}$?
\end{enumerate}
\end{quizquestion}

\begin{solution}
\textbf{(a)} $T\sim t_{n-1}$ under $H_0$.

\textbf{(b)} $h(\mu_0) = P_{H_0}(|T|>t_{n-1;\,1-\alpha/2}) = 2P(t_{n-1}>t_{n-1;\,1-\alpha/2}) = \alpha$. (By symmetry of $t_{n-1}$.)

\textbf{(c)} The $t$-distribution has heavier tails than $\mathcal{N}(0,1)$ (Var$(T_{k})=k/(k-2)>1$). So the upper tail is thicker, meaning you need to go further out to capture the same probability in the tail. Hence the critical value is larger. As $n\to\infty$, $t_{n-1}\to\mathcal{N}(0,1)$ and the values converge.
\end{solution}

\subsection{Challenge Problems}

\begin{quizquestion}{Two-sample problem (preview)}
$X_1,\ldots,X_m\iid\mathcal{N}(\mu_X,\sigma^2)$ and $Y_1,\ldots,Y_n\iid\mathcal{N}(\mu_Y,\sigma^2)$ independent, $\sigma^2$ known. Test $H_0:\mu_X=\mu_Y$ vs $H_1:\mu_X\neq\mu_Y$.
\begin{enumerate}[label=(\alph*)]
  \item What is the distribution of $\overline{X}-\overline{Y}$ under $H_0$?
  \item Write the test statistic and rejection region.
\end{enumerate}
\end{quizquestion}

\begin{solution}
\textbf{(a)} Under $H_0:\mu_X=\mu_Y$: $\overline{X}-\overline{Y}\sim\mathcal{N}(0,\sigma^2/m+\sigma^2/n)=\mathcal{N}(0,\sigma^2(1/m+1/n))$.

\textbf{(b)} $Z=\dfrac{\overline{X}-\overline{Y}}{\sigma\sqrt{1/m+1/n}}\sim\mathcal{N}(0,1)$ under $H_0$.
Reject when $|Z|>z_{1-\alpha/2}$.
\end{solution}

\begin{quizquestion}{Power curve analysis}
$X_1,\ldots,X_{25}\iid\mathcal{N}(\mu,1)$. Bilateral test at $\alpha=0.05$: reject when $|\overline{X}|>0.392$.
\begin{enumerate}[label=(\alph*)]
  \item Verify $0.392 = z_{0.975}/\sqrt{25}$.
  \item Compute power at $\mu=0.5$, $\mu=1$, $\mu=2$.
  \item What happens to power as $|\mu|\to\infty$?
\end{enumerate}
\end{quizquestion}

\begin{solution}
\textbf{(a)} $z_{0.975}/\sqrt{25}=1.96/5=0.392$. \checkmark

\textbf{(b)} Power $=P_\mu(|\overline{X}|>0.392)=P_\mu(|\overline{X}-\mu+\mu|>0.392)$.

At $\mu=0.5$: $Z=5\overline{X}\sim\mathcal{N}(2.5,1)$. Power$=P(|Z|>1.96)=1-[\Phi(1.96-2.5)-\Phi(-1.96-2.5)]=1-\Phi(-0.54)+\Phi(-4.46)\approx 1-0.295+0\approx0.705$.

At $\mu=1$: $Z\sim\mathcal{N}(5,1)$. Power$\approx P(Z>1.96)\approx 1-\Phi(-3.04)\approx 0.999$.

At $\mu=2$: $Z\sim\mathcal{N}(10,1)$. Power$\approx 1$.

\textbf{(c)} Power$\to 1$. A consistent test detects any fixed alternative with probability approaching 1 as $|\mu-\mu_0|\to\infty$.
\end{solution}

\begin{quizquestion}{Connecting Ch.\ 3 and Ch.\ 4}
$X_1,\ldots,X_n\iid\mathcal{N}(\mu,\sigma^2)$ with $\sigma^2$ known. From Chapter 3, the asymptotic normality result gives:
\[
  \frac{\sqrt{n}(\hat\mu-\mu)}{\sigma}\convd\mathcal{N}(0,1).
\]
\begin{enumerate}[label=(\alph*)]
  \item Explain how this directly yields the $Z$-test for $H_0:\mu=\mu_0$.
  \item Why is the $Z$-test exact (not just approximate) for Gaussian data?
  \item For non-Gaussian data with known variance, when can we still use the $Z$-test?
\end{enumerate}
\end{quizquestion}

\begin{solution}
\textbf{(a)} Under $H_0:\mu=\mu_0$, the asymptotic normality gives $\sqrt{n}(\overline{X}-\mu_0)/\sigma\convd\mathcal{N}(0,1)$. Use this as the test statistic; reject when $|Z|>z_{1-\alpha/2}$. This is exactly the $Z$-test.

\textbf{(b)} For Gaussian data, $\overline{X}\sim\mathcal{N}(\mu,\sigma^2/n)$ exactly (not just approximately). So $Z\sim\mathcal{N}(0,1)$ exactly under $H_0$, not just asymptotically. The Type I error is exactly $\alpha$ for any $n$.

\textbf{(c)} By the CLT, for large $n$, $\sqrt{n}(\overline{X}-\mu)/\sigma\approx\mathcal{N}(0,1)$ regardless of the underlying distribution (as long as $\sigma^2<\infty$). So the $Z$-test is approximately valid for any distribution with finite variance when $n$ is large.
\end{solution}

\begin{quizquestion}{Likelihood ratio for Gaussian (bilateral)}
$X_1,\ldots,X_n\iid\mathcal{N}(\mu,\sigma^2)$ with known $\sigma^2$. Write the likelihood ratio $\Lambda=L(\hat\mu)/L(\mu_0)$ and show it simplifies to an increasing function of $M=(\overline{X}-\mu_0)^2/(sigma^2/n)$. This is the derivation from lecture slides 27--28.
\end{quizquestion}

\begin{solution}
\begin{align*}
  L(\mu) &= \prod_{i=1}^n \frac{1}{\sqrt{2\pi\sigma^2}}\exp\!\left(-\frac{(X_i-\mu)^2}{2\sigma^2}\right)
  = (2\pi\sigma^2)^{-n/2}\exp\!\left(-\frac{\sum(X_i-\mu)^2}{2\sigma^2}\right).\\
  \Lambda &= \frac{L(\hat\mu)}{L(\mu_0)}
  = \exp\!\left(-\frac{\sum(X_i-\hat\mu)^2-\sum(X_i-\mu_0)^2}{2\sigma^2}\right).
\end{align*}
Using the identity $\sum(X_i-\mu_0)^2=\sum(X_i-\overline{X})^2+n(\overline{X}-\mu_0)^2$:
\[
  \sum(X_i-\hat\mu)^2-\sum(X_i-\mu_0)^2 = -n(\overline{X}-\mu_0)^2.
\]
So:
\[
  \Lambda = \exp\!\left(\frac{n(\overline{X}-\mu_0)^2}{2\sigma^2}\right) = \exp\!\left(\frac{M}{2}\right),
\]
where $M=n(\overline{X}-\mu_0)^2/\sigma^2$. Since $\exp(M/2)$ is increasing in $M$: $\Lambda>Q\iff M>c$ for some $c$. Under $H_0$: $M=Z^2\sim\chi^2(1)$. \checkmark
\end{solution}

\begin{quizquestion}{Full exam-style problem}
Blood pressure measurements: $X_1,\ldots,X_{20}\iid\mathcal{N}(\mu,25)$ (so $\sigma=5$). A doctor claims the population mean is $\mu_0=120$ mmHg. You observe $\overline{x}=124.3$.

\begin{enumerate}[label=(\alph*)]
  \item Test $H_0:\mu=120$ vs $H_1:\mu\neq 120$ at $\alpha=0.05$. Compute the test statistic, rejection region, and conclusion.
  \item Now test $H_0:\mu\leq 120$ vs $H_1:\mu>120$ at $\alpha=0.05$. Does your conclusion change?
  \item Compute the $p$-value for the bilateral test.
  \item If $\sigma^2$ were unknown and you observed $s=5.2$, redo part (a) using the $t$-test. Use $t_{19;\,0.975}=2.093$.
\end{enumerate}
\end{quizquestion}

\begin{solution}
\textbf{(a) Bilateral test:}
$Z = \dfrac{124.3-120}{5/\sqrt{20}} = \dfrac{4.3}{5/4.472} = \dfrac{4.3}{1.118} \approx 3.85$.

Rejection region: $|Z|>z_{0.975}=1.96$. Since $3.85>1.96$, \textbf{reject $H_0$}. Significant evidence that $\mu\neq 120$.

\textbf{(b) Right-sided test:}
Rejection region: $Z>z_{0.95}=1.645$. Since $3.85>1.645$, \textbf{reject $H_0$}. Same conclusion, but this test is more powerful for detecting $\mu>120$ specifically.

\textbf{(c) $p$-value:}
$p$-value $= P(|Z|>3.85) = 2(1-\Phi(3.85)) \approx 2\times 0.00006 = 0.00012$.

Extremely small --- very strong evidence against $H_0$.

\textbf{(d) $t$-test:}
$T = (124.3-120)/(5.2/\sqrt{20}) = 4.3/(5.2/4.472) = 4.3/1.163 \approx 3.70$.

Rejection region: $|T|>t_{19;\,0.975}=2.093$. Since $3.70>2.093$, \textbf{reject $H_0$}. Same conclusion.

Note: the $t$-statistic (3.70) is slightly smaller than the $Z$-statistic (3.85) because $s=5.2>\sigma=5$, making the denominator slightly larger.
\end{solution}

% ============================================================
\section{Quick Reference Card}
% ============================================================

\begin{tcolorbox}[formulabox]
\textbf{Bilateral $Z$-test} ($\sigma^2$ known): $H_0:\mu=\mu_0$ vs $H_1:\mu\neq\mu_0$
\[
  Z=\frac{\overline{X}-\mu_0}{\sigma/\sqrt{n}}\sim\mathcal{N}(0,1),\quad \delta=\mathbf{1}_{|Z|>z_{1-\alpha/2}},\quad h(\mu_0)=\alpha.
\]
\textbf{Right-sided $Z$-test}: $H_0:\mu\leq\mu_0$ vs $H_1:\mu>\mu_0$
\[
  \delta=\mathbf{1}_{Z>z_{1-\alpha}},\quad h(\mu_0)=\alpha.
\]
\textbf{Left-sided $Z$-test}: $H_0:\mu\geq\mu_0$ vs $H_1:\mu<\mu_0$
\[
  \delta=\mathbf{1}_{Z<-z_{1-\alpha}},\quad h(\mu_0)=\alpha.
\]
\textbf{Bilateral $t$-test} ($\sigma^2$ unknown): $H_0:\mu=\mu_0$ vs $H_1:\mu\neq\mu_0$
\[
  T=\frac{\overline{X}-\mu_0}{S/\sqrt{n}}\sim t_{n-1},\quad \delta=\mathbf{1}_{|T|>t_{n-1;\,1-\alpha/2}}.
\]
\textbf{Type II error (bilateral, $Z$-test)}:
\[
  g(\mu)=\Phi\!\left(z_{1-\alpha/2}-\Delta\right)-\Phi\!\left(-z_{1-\alpha/2}-\Delta\right),\quad \Delta=\frac{(\mu-\mu_0)\sqrt{n}}{\sigma}.
\]
\textbf{$p$-value (bilateral)}: $p=2(1-\Phi(|Z_{\text{obs}}|))=2\Phi(-|Z_{\text{obs}}|)$.\\
\textbf{$p$-value (right-sided)}: $p=1-\Phi(Z_{\text{obs}})=\Phi(-Z_{\text{obs}})$.\\
\textbf{Critical values}: $z_{0.95}=1.645$, $z_{0.975}=1.96$, $z_{0.995}=2.576$.
\end{tcolorbox}

\begin{tcolorbox}[successbox, title={Exam checklist for hypothesis testing problems}]
\begin{enumerate}[leftmargin=2em,itemsep=2pt]
  \item[$\square$] Identify $H_0$ and $H_1$. Which type: simple, unilateral, bilateral?
  \item[$\square$] Is $\sigma^2$ known ($Z$-test) or unknown ($t$-test)?
  \item[$\square$] Write the test statistic ($Z$ or $T$). Identify its distribution under $H_0$.
  \item[$\square$] Determine the rejection region direction (both tails / right / left).
  \item[$\square$] Use $z_{1-\alpha/2}$ for bilateral, $z_{1-\alpha}$ for unilateral. Never mix these up.
  \item[$\square$] Verify $h(\mu_0)=\alpha$ by symmetry argument.
  \item[$\square$] For Type II error: shift the distribution by $\Delta=(\mu-\mu_0)\sqrt{n}/\sigma$ and recompute.
  \item[$\square$] State conclusion: reject $H_0$ / fail to reject $H_0$. Never say ``accept $H_0$''.
\end{enumerate}
\end{tcolorbox}

\end{document}
